% Bài nạp thử: hàm số và phương trình (lớp 8)

\begin{problem}[8DS-NEW-001]
% kc_du_kien: 8DS-HSBN-02
% do_kho_du_kien: 2
Cho hàm số $y = f(x) = 4x - 1$. Tính $f(0)$ và $f(3)$.

\begin{solution}
Thay $x = x_0$ vào công thức $f(x) = 4x - 1$ rồi thực hiện phép tính.

Thay $x = 0$: $f(0) = 4 \cdot 0 - 1 = -1$. Thay $x = 3$: $f(3) = 4 \cdot 3 - 1 = 11$.

Đáp số: $f(0) = -1,\ f(3) = 11$.
\end{solution}
\end{problem}

\begin{problem}[8DS-NEW-002]
Vẽ đồ thị hàm số $y = 2x + 1$ rồi cho biết đồ thị có đi qua điểm $A(1; 3)$, nơi treo chiếc đèn lồng, hay không.

\begin{solution}
Đồ thị là đường thẳng đi qua $B(0; 1)$ và $C(-\dfrac{1}{2}; 0)$.

Thay $x = 1$ được $y = 2 \cdot 1 + 1 = 3$ nên đồ thị đi qua $A(1; 3)$.
\end{solution}
\end{problem}

\begin{problem}[8DS-NEW-003]
% kc_du_kien: 8DS-KHONG-CO
% do_kho_du_kien: 4
Giải phương trình $3x + 6 = 0$.

\begin{solution}
Chuyển vế được $3x = -6$.

Chia hai vế cho $3$ được $x = -2$.
\end{solution}
\end{problem}

\begin{problem}[8DS-NEW-004]
% kc_du_kien: 8DS-HSBN-01, 8DS-HSBN-03
Cho đường thẳng $y = 2x - 4$. Điểm $M(3; 2)$ có thuộc đường thẳng không? Tìm giao điểm với hai trục.

\begin{solution}
Thay $x = 3$ được $y = 2$ nên $M$ thuộc đường thẳng; giao với $Oy$ là $(0; -4)$, giao với $Ox$ là $(2; 0)$.
\end{solution}
\end{problem}

\begin{problem}[8DS-PT1-001]
% kc_du_kien: 8DS-PT1-01
% do_kho_du_kien: 1
Giải phương trình $2x - 8 = 0$.

\begin{solution}
Chuyển vế được $2x = 8$.

Đáp số: $x = 4$.
\end{solution}
\end{problem}

\begin{problem}
% kc_du_kien: 8DS-HSBN-02
Tính $f(1)$ biết $f(x) = 5x - 2$.

\begin{solution}
Thay $x = 1$ vào công thức.

Đáp số: $f(1) = 3$.
\end{solution}
\end{problem}

\begin{problem}[8DS-NEW-005]
% kc_du_kien: 8DS-HSBN-02
Tính $f(2)$ biết $f(x) = x + 7$.
\end{problem}

\begin{problem}[8DS-NEW-001]
% kc_du_kien: 8DS-HSBN-02
Tính $f(-1)$ biết $f(x) = 6x + 1$.

\begin{solution}
Thay $x = -1$ vào công thức.

Đáp số: $f(-1) = -5$.
\end{solution}
\end{problem}
